Drop perpendiculars from the centers to the tangent. Similar right triangles show that the distances from its axis intersection $T$ to the centers have ratio $2:7$. For $0<x<15$, $\dfrac{x}{15-x}=\dfrac{2}{7}$, giving $x=\dfrac{10}{3}$. For $x>15$, the ratio $\dfrac{x}{x-15}>1$ cannot equal $\dfrac{2}{7}$. Thus $\boxed{x=\dfrac{10}{3}}$. The two internal tangents are reflections with the same intersection.
\begin{asy}
texpreamble("\usepackage{amsmath,amssymb}");
size(220);
defaultpen(fontsize(9pt));
pair O=(0,0),C=(15,0),T=(10/3,0);draw(circle(O,2));draw(circle(C,7));draw((-3,0)--(23,0));draw((-2,-4)--(21,13.25));draw((-2,4)--(21,-13.25),dashed);pair U=(1.2,-1.6),V=(10.8,5.6);draw(O--U,dotted);draw(C--V,dotted);dot(O);dot(C);dot(T);label("$O$",O,(-1,1));label("$C$",C,(1,0));label("$T$",T,(1,-1));label("$2$",(O+U)/2,(-1,0));label("$7$",(C+V)/2,(1,0));
\end{asy}
